\section{Конкретно-исторический анализ кризисов и общей политэкономической динамики промышленного цикла}

\subsection{Кризис как конкретно-историческое и социально-экономическое явление}
Кризис перепроизводства в капиталистической экономике представляет собой не просто временное сокращение производства, снижение темпов роста валового продукта или ухудшение отдельных макроэкономических показателей, а исторически обусловленное и внутренне необходимое социально-экономическое явление. Его содержание определяется специфическими противоречиями капиталистического способа производства: общественным характером производственного процесса при частной форме присвоения его результатов, стихийностью рыночного регулирования, ориентацией производства на извлечение прибыли и относительной ограниченностью платежеспособного спроса. Поэтому кризис обнаруживается не только в динамике выпуска, инвестиций, занятости или торговли, но и в массовом обесценении капитала, банкротствах предприятий, росте безработицы, сокращении доходов трудящихся, нарушении кредитных связей и обострении социальных противоречий. Статистические показатели фиксируют лишь отдельные внешние проявления данного процесса и сами по себе не исчерпывают его действительного содержания.

При этом используемые в исследовании статистические данные, как правило, не являются идеальным и непосредственным отражением хозяйственной действительности. На их величину и сопоставимость влияют изменения методик учета, структуры цен, валютных курсов, территориальных границ, отраслевой классификации, полноты регистрации экономической деятельности, а также последствия войн, политических потрясений, природных катастроф, государственных интервенций и иных внешних факторов. Одни и те же количественные изменения могут иметь различную экономическую природу: падение промышленного производства может быть связано как с развертыванием циклического кризиса перепроизводства, так и с военным разрушением производительных сил, административным ограничением выпуска, структурной перестройкой экономики или изменением статистической методологии. Соответственно, механическое отождествление ухудшения отдельных показателей с кризисом общего перепроизводства способно привести к методологически ошибочным выводам.

В силу этого простой эконометрический анализ временных рядов, основанный на поиске статистических минимумов, отклонений от тренда или формальных корреляций, не способен сам по себе дать адекватную характеристику кризисного процесса. Эконометрические методы могут быть полезным вспомогательным инструментом для выявления количественной динамики, однако они не раскрывают причинной структуры явления, его классового содержания и конкретного механизма развития. Формально сходные колебания показателей могут отражать качественно различные процессы, тогда как кризис, уже начавшийся в сфере кредита, торговли или накопления капитала, может еще не получить полного выражения в агрегированных статистических данных. Абстрагирование от исторического контекста и реальных связей воспроизводства превращает количественную модель в описание внешней формы движения, не объясняющее его социально-экономического содержания.

Таким образом, выделение циклических кризисов общего перепроизводства, установление их хронологических границ, а также разграничение периодов и фаз промышленного цикла требуют комплексного конкретно-исторического анализа. Такой анализ должен сочетать исследование динамики промышленного производства, инвестиций, торговли, цен, кредита, прибыли, занятости и банкротств с изучением изменений в структуре капитала, положении основных классов, государственной экономической политике и международной хозяйственной конъюнктуре. Лишь сопоставление количественных данных с качественным анализом исторических условий позволяет отличить собственно циклический кризис капиталистического перепроизводства от нециклических спадов, отраслевых диспропорций и внешне обусловленных хозяйственных потрясений. Следовательно, статистика и эконометрика должны рассматриваться не как самодостаточное основание периодизации цикла, а как элементы более широкого историко-экономического исследования, раскрывающего кризис в единстве его сущности, форм проявления и конкретных условий развития.

\subsection{Типология кризисов капиталистической экономики}
При исследовании промышленного цикла необходимо учитывать не всякое кризисное явление и не любое сокращение производства. Спад может быть вызван войной, структурной перестройкой, изменением государственной политики, нарушением внешней торговли или иными обстоятельствами, непосредственно не связанными с циклическим движением совокупного общественного капитала. Поэтому для периодизации промышленного цикла следует установить масштаб кризиса, его причины и влияние на производство, инвестиции, занятость, торговлю и кредит.

По масштабу распространения различаются кризисы общего и частичного перепроизводства. Кризис общего перепроизводства охватывает капиталистическую экономику в целом и выражается в сокращении производства, инвестиций и занятости, росте товарных запасов, банкротств и безработицы. Кризис частичного перепроизводства, напротив, ограничивается отдельной отраслью или группой отраслей и не обязательно приводит к падению совокупного выпуска. К частичным кризисным явлениям относятся и финансовые кризисы, связанные с падением фондового рынка, нарушением кредитных отношений и обесценением фиктивного капитала, если они не сопровождаются общим сокращением материального производства. Вместе с тем эти формы тесно взаимосвязаны: общий кризис нередко начинается в отдельных отраслях, крупнейших корпорациях, банковской системе или на рынке ценных бумаг, а затем распространяется на всю экономику.

По характеру причин кризисы общего перепроизводства подразделяются на циклические и промежуточные. Циклические кризисы обусловлены внутренними противоречиями капиталистического воспроизводства и периодичностью обновления основного капитала. Расширение производственных мощностей в условиях стихийности и конкуренции постепенно приводит к перенакоплению капитала, перепроизводству товаров и невозможности их прибыльной реализации. Промежуточные кризисы также могут охватывать значительную часть экономики, однако их непосредственный импульс связан преимущественно с внешними факторами: войнами, энергетическими и сырьевыми шоками, эмбарго, резкими изменениями мировой торговли или государственной политики. Они обостряют внутренние противоречия капитализма, но не обязательно завершают полный промышленный цикл.

Таким образом, при анализе промышленного цикла необходимо разграничивать частичные, финансовые и общие кризисы, а внутри последних — циклические и промежуточные. Основными критериями такого разграничения выступают масштаб спада, его воздействие на воспроизводство совокупного общественного капитала, связь с обновлением основного капитала и конкретно-исторические причины возникновения. Это позволяет отделить собственно циклический кризис общего перепроизводства от локальных, финансовых и внешне обусловленных потрясений.

\subsection{Периоды промышленного цикла}
Продолжительность периода промышленного цикла определяется временным интервалом между двумя последовательными циклическими кризисами общего перепроизводства. Исходным пунктом такого исчисления является кризис, поскольку именно в нем с наибольшей полнотой обнаруживаются накопившиеся противоречия капиталистического воспроизводства и одновременно формируются материальные предпосылки последующего цикла. Следовательно, период промышленного цикла охватывает движение капиталистической экономики от одного циклического кризиса до следующего, включая сменяющие друг друга состояния депрессии, оживления и подъема. При этом продолжительность цикла не следует смешивать ни с продолжительностью собственно кризисной фазы, ни с длительностью статистически фиксируемого сокращения производства: кризис образует границу между циклами, тогда как сам цикл представляет собой более продолжительный процесс воспроизводства и накопления общественного капитала.

Для количественного определения продолжительности каждого цикла необходимо последовательно сопоставить даты начала двух смежных циклических кризисов общего перепроизводства. Если обозначить дату начала предыдущего кризиса через $t_i$, а дату начала следующего – через $t_{i+1}$, то продолжительность соответствующего периода промышленного цикла составит $T_i=t_{i+1}-t_i$. Средняя продолжительность определяется как среднее арифметическое всех завершенных межкризисных интервалов. Принципиально важно применять единую систему датировки: если границей одного цикла считается месяц начала общего промышленного спада, то аналогичный критерий должен использоваться и для всех остальных кризисов. Иначе различия в методике определения кризисных дат могут искусственно увеличивать или сокращать измеряемые интервалы.

Согласно принятой в настоящей работе периодизации циклических кризисов общего перепроизводства, средняя продолжительность периода промышленного цикла в экономике США с 1919 по 2026 г. составляет 120,1 месяца, то есть приблизительно десять лет. Тем самым полученный результат практически идеально соответствует десятилетнему периоду. Вместе с тем отдельные циклы существенно отклоняются от среднего значения: максимальная продолжительность межкризисного интервала достигает 148 месяцев, или 12 лет и 4 месяцев, тогда как минимальная составляет 81 месяц, или 6 лет и 9 месяцев. Размах между максимальным и минимальным значениями равен 67 месяцам, что свидетельствует о том, что промышленный цикл не представляет собой строго механический и хронологически неизменный процесс.

Различия в продолжительности отдельных циклов обусловлены конкретно-историческими условиями капиталистического воспроизводства. На сроки развертывания кризиса воздействуют темпы накопления капитала, масштабы расширения производственных мощностей, уровень загрузки оборудования, состояние кредита, динамика нормы прибыли, государственная экономическая политика и положение страны в мировом хозяйстве. Существенное влияние оказывают также войны, энергетические и сырьевые потрясения, структурные сдвиги, изменения международной торговли и другие внешние обстоятельства. Эти факторы могут ускорить созревание кризиса либо, напротив, временно отсрочить его наступление. Поэтому десятилетняя продолжительность должна рассматриваться не как жесткая календарная закономерность, а как средняя историческая форма проявления внутренних закономерностей движения капитала.

Полученная средняя продолжительность цикла – около десяти лет – поэтому согласуется с порядком величины срока воспроизводства и физического износа активной части основного капитала в промышленном производстве. Вместе с тем данное соответствие не следует понимать как тождество между сроком службы каждой отдельной машины и продолжительностью всего промышленного цикла. Различные элементы основного капитала имеют неодинаковые сроки эксплуатации: оборудование изнашивается быстрее, чем здания и сооружения, а его фактическое выбытие определяется не только физическим разрушением, но и моральным устареванием, изменением технологии и условиями конкурентной борьбы. В статистической методологии износ основного капитала также рассматривается как результат совокупного воздействия физического изнашивания, старения, повреждений и устаревания.  Речь, следовательно, идет о среднем общественном сроке воспроизводства наиболее активной части промышленного капитала, а не о едином техническом нормативе для всех средств производства.

Таким образом, установленные межкризисные интервалы подтверждают наличие устойчивого, хотя и не абсолютно регулярного, промышленного цикла. Среднее значение в 120,1 месяца выражает общую десятилетнюю периодичность движения капиталистического воспроизводства, тогда как отклонения от него отражают воздействие конкретных исторических условий. Единство средней закономерности и изменчивости ее отдельных проявлений соответствует диалектическому характеру промышленного цикла: его материальная основа воспроизводится через оборот основного капитала, но конкретная продолжительность каждого периода определяется всей совокупностью экономических и исторических обстоятельств. $\dagger$ (Приложение \ref{app:emp}, стр. \pageref{fig:indpro}, \pageref{fig:phases-duration})

\subsection{Фазы промышленного цикла}
Промышленный цикл представляет собой закономерную последовательность качественно различных состояний капиталистического воспроизводства. В классическом виде он включает четыре основные фазы: кризис, депрессию, оживление и подъем. Каждая из них характеризуется определенным направлением движения промышленного производства, инвестиций, занятости, цен, прибыли и кредита. Таким образом, фаза кризиса характеризуется общим спадом промышленного производства, фаза депрессии – его замедлением дальнейшим застоем производства, фаза оживления – началом роста, а фаза подъема – превышением докризисного максимума. Вместе с тем границы между фазами не являются абсолютно резкими: процессы, характерные для предшествующей фазы, могут некоторое время сохраняться после начала следующей, а изменение отдельных показателей происходит не одновременно. Поэтому фазу цикла следует определять не по динамике одного статистического ряда, а на основании совокупности количественных и качественных признаков, отражающих состояние воспроизводства общественного капитала. Следовательно, стандартная последовательность \textquotedbl{}кризис $\to$ депрессия $\to$ оживление $\to$ подъем\textquotedbl{}, отражающая внутреннюю логику движения промышленного цикла, в реальной истории капиталистической экономики может претерпевать некоторые изменения.

Первой формой такого отклонения являются промежуточные кризисы, которые чаще всего возникают в ходе фазы подъема. Они сопровождаются общим сокращением производства, инвестиций и занятости, однако непосредственно не вытекают из завершения очередного цикла массового воспроизводства основного капитала. Их возникновение в значительной степени связано с внешними или экстраординарными обстоятельствами, которые воздействуют на уже существующие противоречия капиталистической экономики. Промежуточный кризис может нарушить поступательное развитие подъема, но его дальнейшие последствия зависят от глубины спада и общего состояния процесса накопления. После промежуточного кризиса возможны два основных варианта движения. В первом случае кризис вызывает временное сокращение производства, но не завершает фазу подъема. Подобная динамика наблюдалась после промежуточного кризиса 1973 г.: вызванное внешним воздействием обострение хозяйственных противоречий прервало подъем, но не стало окончательной границей соответствующего промышленного цикла. Во втором случае промежуточный кризис может завершить подъем, если его воздействие совпадает с накоплением значительных внутренних диспропорций и приводит к глубокому изменению условий общественного воспроизводства. Так, кризис 1945 г., связанный с окончанием Второй мировой войны, прекращением значительной части военных заказов, демобилизацией и конверсией промышленности, не просто временно прервал расширение производства, но завершил предшествующую фазу подъема

Другой разновидностью отклонения от стандартной последовательности является кратковременное промежуточное оживление внутри затяжного кризиса общего перепроизводства. После первоначального сокращения производства государство может перейти к стимулирующей денежно-кредитной и бюджетной политике либо ослабить ранее проводившиеся ограничительные меры. Снижение процентных ставок, расширение государственных расходов, налоговые послабления, поддержка кредитной системы и увеличение государственных заказов способны временно повысить совокупный спрос, стабилизировать инвестиции и вызвать кратковременное увеличение промышленного выпуска. Однако поскольку накопившееся перенакопление капитала, долговые противоречия и отраслевые диспропорции не были в достаточной степени разрешены, первоначальное восстановление оказывается неустойчивым. После кратковременного роста производство вновь начинает сокращаться, и кризис продолжается. Подобная прерывистая динамика была характерна для кризисных периодов 1957–1961 и 1980–1982 гг. В обоих случаях временное улучшение конъюнктуры, связанное в том числе с изменением направленности государственной политики, не устранило материальных причин кризиса и сменилось новым падением производства.

Средняя продолжительность фаз промышленного цикла за исследуемый период в экономике США составляет 16,6 месяцев кризиса (13,4 с учетом промежуточных оживлений), 3,3 месяца депрессии, 23,0 месяца оживления, и 79,5 месяцев подъема (64,1 с учетом промежуточных кризисов). Таким образом, наиболее короткой из фаз является депрессия, а наиболее продолжительной – подъем. При этом стоит отметить, что за исследуемый период продолжительность фазы депрессии сократилась с 6-8 месяцев до 1-2 месяцев, что можно связать с развитием стимулирующих мер антикризисной политики.

Таким образом, стандартная последовательность фаз промышленного цикла выражает закономерное движение от кризисного сокращения через стабилизацию и восстановление к новому расширению производства. Вместе с тем конкретная история цикла включает промежуточные кризисы и временные оживления, которые могут нарушать внешнюю последовательность фаз, не отменяя ее внутренней логики. Полученные показатели продолжительности подтверждают, что наиболее краткой фазой является депрессия, а наиболее продолжительной – подъем, в ходе которого постепенно воспроизводятся материальные условия нового кризиса общего перепроизводства. $\dagger$ (Приложение \ref{app:emp}, стр. \pageref{fig:indpro}, \pageref{fig:phases-duration})

\subsection{Особенности текущего промышленного цикла}
Текущий промышленный цикл в экономике США, начавшийся после кризисного сокращения производства 2020 г., имеет ряд существенных особенностей, отличающих его как от классической схемы промышленного цикла, так и от большинства предшествующих циклов исследуемого периода. К их числу относятся чрезвычайно быстрое прекращение первоначального кризисного падения, продолжительное восстановление промышленного производства, отсутствие устойчивого превышения максимума предыдущего цикла, длительная стагнация физического объема выпуска и возрастающая роль финансового накопления, связанного с развитием искусственного интеллекта. Совокупность этих обстоятельств свидетельствует о противоречивом характере текущего цикла: формальное расширение хозяйственной активности сочетается со слабостью материально-производственной основы накопления.

Согласно официальной хронологии Национального бюро экономических исследований США, последний максимум общей экономической активности был достигнут в феврале 2020 г., а минимум – в апреле 2020 года. Поэтому май 2020 г. рассматривается NBER как первый месяц последующего расширения. При этом необходимо учитывать, что датировка рецессий NBER относится к фактическому состоянию кризиса общего перепроизводства и основывается на совокупности показателей доходов, занятости, потребления, торговли и производства, что не является тождественным периодизации промышленного цикла исключительно исходя из динамики промышленного производства.

Формально экономика находится в фазе подъема относительно последнего кризиса 2020 (если рассматривать момент общего кризиса исходя из датировок NBER), однако по факту докризисный максимум 2018 года в 104 пункта так и не был преодолен. Более того, американская экономика едва ли сумела преодолеть докризисный уровень 2007 года всего лишь два раза, в 2014 и 2018 годах соответственно, и на текущий момент значение индекса промышленного производства по-прежнему уступает максимуму 2007 года. Помимо этого, в предпоследнем и ткущем промышленных циклах фаза оживления, которая в среднем за исследуемый период составляет около 23 месяцев (почти два года), заняла 59 и 60 месяцев соответственно (почти пять лет, что составляет половину среднего периода всего промышленного цикла). Таким образом, физический объем промышленного производства в экономике США фактически не растет с 2007 года, а с 2022 года по сути находится в состоянии стагнации: в апреле 2022 г. индекс промышленного производства составлял 101,44 пункта, тогда как в течение 2023–2025 гг. он преимущественно колебался в сравнительно узком диапазоне около 99–102 пунктов, а к маю 2026 года поднялся до 102,65 пункта, то есть оказался лишь приблизительно на 1,2 пункта выше уровня апреля 2022.

Исходя из выявленной в настоящей работе исторической периодичности, следующий циклический кризис общего перепроизводства в США при прочих равных условиях может развернуться в период с 2028 по 2032 г. Данный интервал соответствует наступлению кризиса приблизительно через восемь-двенадцать лет после 2020 года и находится в пределах исторически наблюдавшейся продолжительности промышленного цикла. Однако по состоянию на середину 2026 года более точное определение даты следующего кризиса представляется преждевременным, поскольку большинство статистических показателей еще не образует единой совокупности признаков, однозначно свидетельствующей о непосредственном приближении кризиса общего перепроизводства. Циклическая периодичность не действует как механический календарный закон. Наступление кризиса зависит от темпов накопления и обновления основного капитала, движения нормы прибыли, кредитной экспансии, долговой нагрузки, динамики инвестиций, товарных запасов, загрузки мощностей и международной конъюнктуры. Государственная политика способна на определенное время отсрочить открытое проявление кризиса, тогда как финансовое потрясение, война, сырьевой шок или нарушение мировой торговли могут ускорить его наступление. Поэтому интервал 2028–2032 гг. следует рассматривать как авторский сценарный прогноз, основанный на средней и предельной продолжительности прошлых циклов.

Одним из возможных непосредственных поводов для нового кризиса может стать обесценение капитала, вложенного в сферу искусственного интеллекта. В теории искусственный интеллект способен повышать производительность труда, изменять организацию производства и создавать новые средства производства; однако на практике действительный технологический потенциал в рамках капиталистической экономики подменяется спекулятивной формой капитализации технологий и образованием финансового пузыря, основанного на завышенных ожиданиях будущей прибыли. К началу 2026 г. накопление капитала в данном секторе приобрело значительные масштабы. По оценке ОЭСР, в 2025 г. на компании, связанные с искусственным интеллектом, приходился 61\% стоимости всех мировых венчурных инвестиций, тогда как в 2022 г. их доля составляла 30\%. Около 75\% мировой стоимости венчурных сделок с компаниями сферы искусственного интеллекта приходилось на США. При этом 73\% совокупной стоимости подобных инвестиций составляли крупные сделки объемом более 100 млн долларов, что указывает на высокую концентрацию капитала.

Когда доходы от применения искусственного интеллекта окажутся ниже ожидаемых, начавшееся обесценение акций может распространиться далеко за пределы технологического сектора. Снижение рыночной капитализации ухудшит возможности компаний привлекать финансирование, вызовет пересмотр инвестиционных программ, сократит заказы на полупроводники, вычислительное оборудование, строительство центров обработки данных и энергетическую инфраструктуру. Одновременно убытки банков, инвестиционных фондов и частных кредиторов могут привести к общему ужесточению кредита. Посредством взаимосвязи фиктивного и финансового капитала частичный кризис в сфере искусственного интеллекта способен перерасти в кризис общего перепроизводства. Однако финансовый пузырь, разрушение которого послужит непосредственным толчком к спаду, выступает исключительно в качестве повода, а не первопричины циклического кризиса, поскольку обесценение фиктивного капитала лишь открыто проявляет противоречия, ранее накопившиеся в процессе реального воспроизводства: перенакопление капитала, несоответствие между расширением производственных мощностей и платежеспособным спросом, рост задолженности и затруднение прибыльного применения капитала; в то время как финансовый кризис выступает лишь механизмом, посредством которого эти противоречия переходят из скрытой формы в открытую. $\dagger$ (Приложение \ref{app:emp}, стр. \pageref{fig:indpro}, \pageref{fig:phases-duration})